\FloatBarrier
\section{Conclusion}
\label{sec:conclusion}

We presented \model{}, which decides NLI from the coverage of each hypothesis word by the premise and
from conflicts between them, using contextual token states aligned through typed lexical relations. The
structured decision is added to the logits of the standard classifier, at a cost of about 2M parameters.

Across English and Portuguese, with large encoders, 10 seeds and paired tests, \model{} matches the
\cls{} baseline in distribution and improves on it out of distribution, by up to 8.6 MCC points in English
and 2.3 in Portuguese, with lower variance across seeds. In English it also outperforms three
knowledge-graph baselines, which fall below \cls{} out of distribution.

In English, the ablation locates the effect in the coverage term and the logical aggregation; in
Portuguese, with a larger encoder, no single component is needed; the conflict term,
the lexicon and its source contribute little. The AX-b analysis shows where the rule helps, on argument
structure and logical operators, and where it fails, on restrictivity, universal quantifiers and
downward monotonicity.

Two directions follow. Extending coverage with the monotonicity of the context would address its main
failure mode, and the conflict term needs either a different training signal or removal. Cross-lingual
transfer is the weak point: the structured branch learned in English does not transfer to Portuguese,
and from Portuguese to English its gains shrink to ConjNLI.
\todo{revise when experiment 0014 finishes (Portuguese and cross-lingual baselines).}
